\section{Related work}
\label{sec:related}

\paragraph{Run-to-run variance.} That RL results move under reseeding is old news in deep RL
\citep{henderson2018matters} and documented again for language-model reasoning
\citep{hochlehnert2025sober}. Those papers establish the size of the problem by replication, which
is the only instrument available without a model of where the variance comes from.
\citet{agarwal2021precipice} prescribe how the resulting samples should be summarised. All of it
presumes the runs exist. Practice for benchmark error bars treats the evaluation set as the only
source of randomness \citep{miller2024adding}, which is the right interval for a fixed model and
the wrong one for a training procedure; \citet{bouthillier2021accounting} argue the pipeline's own
randomness belongs in the interval and estimate it by replication. We estimate the same quantity
without replicating, and our real-model measurements subtract the binomial term explicitly.

\paragraph{Deciding when to stop repeating.} Stopping a comparison once one arm is convincingly
ahead is sequential analysis \citep{wald1945sequential}, and its machine-learning form is racing
\citep{maron1993hoeffding, evendar2006action}, sharpened by using an empirical variance in the
bound \citep{audibert2009exploration} and scaled to hyperparameter search by allocating a budget
across candidates \citep{li2017hyperband}. Section~\ref{sec:sizing} uses the simplest rule of that
family deliberately: the contribution is not a better stopping rule but the quantity the rule
consumes. Every method there takes its variance from samples already drawn, which is sound when a
sample is a validation example and unaffordable when a sample is a training run. At one run per arm
those bounds are vacuous, and the experiment shows that the resulting floor, not the shape of the
bound, is what a comparison pays.

\paragraph{Stochastic dynamics of training.} Treating stochastic optimisation as a process whose
covariance is set by gradient noise against curvature is standard \citep{mandt2017sgd}, and
\eqref{eq:lyapunov} is not new as mathematics. What we add is the RLVR-specific measurement of
$Q_t$ through the two-level decomposition, the finite-time rather than stationary treatment, and
the observation that the filter, not the size of the noise, is what governs reproducibility here.
\citet{wang2026linear} reach a filtering picture from the opposite side, reporting near-linear
evolution of weights and log-probabilities; their object is one run's trajectory, ours the
covariance between runs differing only in seed.

\paragraph{Adjoint sensitivity.} Carrying a scalar backwards through a trajectory to avoid forming
a large intermediate is the adjoint method \citep{pontryagin1962mathematical, chen2018neural}.
Theorem~\ref{thm:adjoint} applies it to a covariance rather than a loss, and
Proposition~\ref{prop:potential} is what makes it cheap here by removing the transpose.

\paragraph{Entropy collapse.} That RLVR drives policy entropy down sharply is documented
\citep{cui2025entropy} and treated as an optimisation pathology. It has a second consequence that
follows from \eqref{eq:adjoint}: a count-based advantage is identically zero on a unanimous group,
so a collapsed policy injects $\Sigma_t = 0$, stops accumulating seed uncertainty, and any
measurement of its reproducibility from that point is measuring nothing.

\paragraph{Batch size and gradient noise.} The gradient noise scale predicts the batch size beyond
which data parallelism stops buying speed \citep{mccandlish2018empirical, shallue2019measuring}.
Both concern a single level of sampling. RLVR samples prompts and then responses within each
prompt, and the resulting two-level scale has no counterpart there; the decomposition itself is the
classical two-stage sampling result \citep{cochran1977sampling}.

\paragraph{Allocating rollouts.} Recent methods vary responses per prompt within a batch from pass
rates or gradient variance \citep{nguyen2026allocation}, or adapt batch size from a divergence
signal \citep{park2026batchscaling}. Those distribute responses \emph{within} one update, across
its prompts. The budget at issue here is across independent training runs of the same recipe, which
is a level those methods do not address and where the reproducibility cost is actually incurred.
